\subsection{Exercices}

\targetExo{fsec2_ex1}
\begin{exercicebox}[Exercice 1 : Arbitrage sur la parité Call-Put]
Sur le marché, on observe les paramètres suivants : Spot $S_0 = 100$, Strike $K = 100$, Échéance $T = 1$ an, $r = 5\%$ (continu). Le Call cote $C = 10$ et le Put cote $P = 8$. L'action ne verse pas de dividendes.
Montrez que ces prix violent la relation de parité. Définissez précisément le portefeuille à constituer aujourd'hui pour réaliser un profit d'arbitrage, et prouvez que le risque final est nul.
\navToCorr{fsec2_ex1}
\end{exercicebox}

\targetExo{fsec2_ex2}
\begin{exercicebox}[Exercice 2 : Borne inférieure stricte du Call Européen]
En utilisant un portefeuille composé d'une option d'achat (Call européen), de l'action sous-jacente et d'un actif sans risque, démontrez rigoureusement que le prix du Call $C_0$ satisfait toujours l'inégalité :
\begin{equation*}
    C_0 \geq \max(0, S_0 - K e^{-rT})
\end{equation*}
\navToCorr{fsec2_ex2}
\end{exercicebox}

\targetExo{fsec2_ex3}
\begin{exercicebox}[Exercice 3 : Borne inférieure stricte du Put Européen]
Par un raisonnement similaire à celui de l'exercice précédent, ou en utilisant la relation de parité Call-Put, démontrez rigoureusement que le prix d'un Put européen $P_0$ satisfait toujours :
\begin{equation*}
    P_0 \geq \max(0, K e^{-rT} - S_0)
\end{equation*}
\navToCorr{fsec2_ex3}
\end{exercicebox}

\targetExo{fsec2_ex4}
\begin{exercicebox}[Exercice 4 : Inégalité stricte des écarts de Strike]
Soit deux options d'achat européennes (Calls) de même maturité $T$ et de prix d'exercice respectifs $K_1$ et $K_2$ avec $K_1 < K_2$. Démontrez, en construisant un portefeuille combinant les deux Calls (un "Bull Spread"), que la différence de prix entre ces deux options est bornée supérieurement par la différence actualisée de leurs strikes :
\begin{equation*}
    C_0(K_1) - C_0(K_2) \leq (K_2 - K_1)e^{-rT}
\end{equation*}
\navToCorr{fsec2_ex4}
\end{exercicebox}